% -------------------------------------------------------------
%   Sprints
% -------------------------------------------------------------

\noindent
This chapter presents the work carried out during each of the project's
six sprints. Each sprint's goal, team composition, key Daily Standup updates,
Sprint Review outcome, and Sprint Retrospective conclusions are reported.

% =============================================================
\section{Sprint 1}
\label{sec:sprint1}
% =============================================================

The first sprint focused on establishing the foundation of the project. The
sprint goal was \enquote{Create design tools for all of the modules and
establish the project infrastructure}, which consisted of the Noun--Verb
analysis, CRC cards and UML diagrams for each of the 5 modules (Asset
Manager, Source Data Manager, Result Data Manager, Optimizer and the
Data Visualizer). Alongside was established the project infrastructure
work: linking Jira with GitHub, creating the initial Avalonia project,
converting the XLSX file with data into a CSV file and storing it in the
assets folder.

Given the varying levels of experience and skillset within the group, the team was divided
into three pairs:
\begin{itemize}
    \item Maria Tsioni and Zoja Angelovicova
    \item Pavel Seliavko and Zoe Elea Schiwy
    \item Dominik Bezo and Daniel Falat
\end{itemize}

Dominik Bezo acted as the Scrum Master, Pavel Seliavko acted as the
Meeting Coordinator and after discussion with the supervisor, Radoslaw
Marek Chylak, the role of the Product Owner was assigned to Radoslaw
Marek Chylak.

The first Daily Standup revealed that additional tasks could be incorporated into the sprint.
After consulting with the supervisor a Figma prototype was added to the Sprint Backlog.
During the Sprint Review, all design tool tasks had the Noun--Verb analysis approved and
marked as done. However, several CRC cards and UML diagrams required further
improvements and were consequently moved back to the Product Backlog (see Appendix~\ref{sec:appendix-backlogs} Fig.~\ref{fig:sprint1-backlog-end}).

The primary challenges were addressed during the Sprint Retrospective:
\begin{itemize}
    \item Dismissiveness and resistance to different ideas,
    \item Low engagement levels during online meetings,
    \item Lack of timely interpersonal feedback among team members,
    \item Underestimation of the workload.
\end{itemize}

These issues were addressed by adopting a more receptive approach, requiring cameras to be on during online meetings, and encouraging immediate escalation of issues rather than deferring them to the Sprint Retrospective. Additionally, a custom story point system was introduced, rating each task by both difficulty and importance on a scale of 1 to 5, with the aim of improving collaboration and workload estimation in future sprints.

% =============================================================
\section{Sprint 2}
\label{sec:sprint2}
% =============================================================

Following a clearer understanding of the workflow and the refinement of certain design tools, Sprint 2 focused on the actual implementation of the first two of the modules which were the Asset Manager and the Source Data Manager. Building on the Figma prototype, completed in Sprint 1, the first version of the Graphical User Interface was also built during this sprint. Each task was also further divided into smaller subtasks (see Appendix~\ref{sec:appendix-backlogs} Fig.~\ref{fig:sprint2-backlog-start}).

Pavel Seliavko assumed the role of the Scrum Master while also keeping the role of the Meeting Coordinator. The original three pairs from Sprint 1 remained unchanged. Using the agreed story point system, each team independently rated every ticket on both importance and difficulty (e.g. the Source Data Manager ticket received an importance score of 5 and a difficulty score of 4).

Early on in the sprint, during the first Daily Standup meeting, it was recognized that the initial plan, which was focusing on the Data Visualizer, was not feasible, therefore the task was replaced with the implementation of the Source Data Manager. This adjustment highlighted that the initial planning had underestimated the complexity of the entire implementation.

During the Sprint Retrospective, it was identified that one of the pair-programming teams experienced difficulty finding common working hours due to conflicting schedules, which hindered the progression of their assigned tickets. To address this, pair-programming assignments were reassigned in the following sprint, deliberately pairing team members with more compatible schedules.

Sprint 2 successfully delivered the Asset Manager, the Source Data Manager and the first version of Graphical User Interface (see Appendix~\ref{sec:appendix-backlogs} Fig.~\ref{fig:sprint2-backlog-end}). The initial burndown chart appeared incorrect because Jira did not register the completed subtasks until the whole tickets were moved to \enquote{DONE} category.

\begin{figure}[htbp]
    \centering
    \includegraphics[width=0.95\textwidth]{sprint2-burndown-jira.png}
    \caption{Sprint 2 burndown chart.}
    \label{fig:sprint2-burnup}
\end{figure}

% =============================================================
\section{Sprint 3}
\label{sec:sprint3}
% =============================================================

Sprint 3 focused on the remaining modules which were the Optimizer, the Result Data Manager and the Data Visualizer. Zoja Angelovicova acted as the Scrum Master and Dominik Bezo as the Meeting Coordinator, while Pavel Seliavko transitioned into the role of internal Product Owner. To address the mid-sprint creation and replacement of the tickets, a problem identified in Sprint 2, the planning style was changed. Each pair created its own detailed subtickets for the assigned module before starting with the development. To address the burndown chart inaccuracies encountered in Sprint 2, each subtask was logged as an individual ticket to ensure accurate progress tracking in Jira (see Appendix~\ref{sec:appendix-backlogs} Fig.~\ref{fig:sprint3-backlog-start}).

Additionally, pair-programming assignments were reorganised to ensure greater schedule compatibility among team members:
\begin{itemize}
    \item Maria Tsioni and Dominik Bezo
    \item Zoja Angelovicova and Daniel Falat
    \item Pavel Seliavko and Zoe Elea Schiwy
\end{itemize}

Sprint 3 successfully delivered a functioning Optimizer, Result Data Manager, and Data Visualizer. However, integration errors encountered during the merging of parallel branches required the sprint deadline to be extended from April 16th to April 19th. As a result, several tickets remained incomplete at the originally planned end of the sprint (see Appendix~\ref{sec:appendix-backlogs} Fig.~\ref{fig:sprint3-backlog-16apr}). The conflicts were resolved by a full team merge. All of the branches were merged into a single branch \texttt{SCRUM-81-Version-1}. Furthermore several tickets had to be renamed to reflect what had actually been implemented.

The Retrospective used a \textsc{Start} / \textsc{Stop} / \textsc{Continue} board with silent writing and round-robin sharing methods (see Appendix~\ref{sec:appendix-backlogs} Table~\ref{tab:sprint3-retro}).

% =============================================================
\section{Sprint 4}
\label{sec:sprint4}
% =============================================================

The primary objective of the Sprint 4 was to finalise and test the product. The nature of the tasks in this sprint was better suited to two groups of three members rather than three pairs. This reorganisation was carried out in accordance with the recommendation made during the previous Sprint Retrospective:
\begin{itemize}
    \item \textbf{Team Slovaks} -- Daniel Falat, Dominik Bezo and Zoja Angelovicova
    \item \textbf{Team International} -- Maria Tsioni, Pavel Seliavko and Zoe Elea Schiwy
\end{itemize}

Daniel Falat took the Scrum Master role and Dominik Bezo took the role of the Meeting Coordinator. Planned work included switching between scenario 1 and scenario 2, refactoring of the Optimizer, removing redundancy in the TimeSeries, implementing maintenance selection, updating the Result Data Manager headers, applying the UML feedback from the lecturer, and writing unit, functional and integration tests for every module. Time estimation and the time actually spent was also added to the tickets for the first time.

Because the focus was heavily directed at making the technical product fully working the tickets regarding report-writing originally planned for this sprint had to be returned to the backlog at the end of the sprint (see Appendix~\ref{sec:appendix-backlogs} Fig.~\ref{fig:sprint4-completed} and Fig.~\ref{fig:sprint4-incomplete}).

The primary issue of the sprint surfaced during the Task Review: the two teams had independently produced very similar solutions for different tasks (SCRUM 117 \& 118 and SCRUM 128 \& 131), creating duplicate code. The conflict was resolved through a discussion between the whole team where it was decided that for documentation reasons, the best resolution was to keep the duplicate implementation while only keeping one solution active. By the Sprint Review on the 1st of May the Minimum Viable Product was complete and verified through testing. Report writing had to be transferred to Sprint 5. The team determined that a comprehensive Start / Stop / Continue retrospective was not warranted this sprint, as collaboration within both teams had functioned effectively.

% =============================================================
\section{Sprint 5}
\label{sec:sprint5}
% =============================================================

Sprint 5 focused mainly on the division of the report tasks, report-writing and refactoring of the code. Because of the Math 2 exam on the 8th of May, a thorough Sprint Planning meeting was held before that on the 3rd of May. The sprint was formally started on the 9th of May, with the end date moved to the 17th of May. Maria assumed the role of Scrum Master and Zoja the Meeting Coordinator, Pavel continued to be the Product Owner. The team composition established in Sprint 4 was retained, maintaining the same two groups of three members (see Appendix~\ref{sec:appendix-backlogs} Fig.~\ref{fig:sprint5-backlog-start}).

On the development side, Team Slovaks was responsible for ensuring the Data Visualizer called Optimizer only once with a clearer separation of concerns, and Team International took the refactoring of the Unit models to use more abstraction/inheritance and fewer interfaces. For the report, every chapter was assigned to an individual author and reviewer (see Table~\ref{tab:report-responsibilities} in Appendix~\ref{sec:report-responsibilities}), except for Chapter 5 which is the Implementation. This chapter is jointly authored by both teams. A dedicated meeting was scheduled for the 16th of May to produce the final Class, Activity and Sequence diagrams used in the report.

During the Sprint Review and Retrospective, no major issues were identified. However, the report-writing tasks had to be carried over to the next sprint, with the tickets already being in progress.

% =============================================================
\section{Sprint 6}
\label{sec:sprint6}
% =============================================================

Sprint 6 was the final sprint of the project, scheduled from the 18th of May to the 1st of June. It was set to be longer than the previous ones because it had to cover the last round of code changes, writing of the remaining report chapters and the full review cycle before the submission deadline.

The two teams from Sprint 4 and 5 were kept, but the team structure was only used for the shared Implementation chapter. The rest of the report writing was assigned per chapter, with a dedicated reviewer for each one. Zoe Elea Schiwy assumed the role of the Scrum Master, Daniel Falat the role of Meeting Coordinator, and Pavel Seliavko remained the Product Owner.

The sprint goal was \enquote{Finish the Report and change the code based on the feedback from Danfoss}. The code-related part of the goal followed from the final demonstration with Danfoss representatives, during which three concrete points were raised: the ratio used in the Optimizer had to be adjusted, and the legend and colours in the Data Visualizer had to be corrected. These were tracked as SCRUM-170, SCRUM-171 and SCRUM-173.

Additional feedback was received from lecturer Amir Ghomashi and distributed across the existing report-writing tickets rather than turned into a separate ceremony. No major conflicts arose during Sprint 6, and the Sprint Review and Retrospective were planned for the day before the submission deadline.